% part: intuitionistic-logic
% chapter: semantics
% section: introduction

\documentclass[../../../include/open-logic-section]{subfiles}

\begin{document}

\olfileid{int}{sem}{int}

\olsection{ভূমিকা}

অর্থবিজ্ঞান ছাড়া কোনো যুক্তিবিদ্যার বিবরণই সন্তোষজনক হয় না;
স্বজ্ঞাবাদী যুক্তিবিদ্যাও এর ব্যতিক্রম নয়। ধ্রুপদি যুক্তিবিদ্যায়
!!{valuation}s-ভিত্তিক অর্থবিজ্ঞান প্রামাণ্য হলেও স্বজ্ঞাবাদী
যুক্তিবিদ্যার জন্য একাধিক প্রতিদ্বন্দ্বী অর্থবিজ্ঞান আছে।
সংযোজকগুলির অর্থের স্বজ্ঞাবাদীর কাছে গ্রহণযোগ্য বিবরণ দেওয়ার
বিচারে এগুলির কোনোটিই সম্পূর্ণ সন্তোষজনক নয়।

মোডাল যুক্তিবিদ্যার অর্থবিজ্ঞানের মতো !!{relational model}s-ভিত্তিক
অর্থবিজ্ঞান সম্ভবত সবচেয়ে প্রচলিত। এখানে !!{propositional variable}s-কে
জগতের সঙ্গে যুক্ত করা হয়, আর জগৎগুলি একটি
অভিগম্যতা-সম্পর্কে আবদ্ধ থাকে। সেই সম্পর্ক সর্বদা একটি আংশিক ক্রম:
অর্থাৎ প্রতিবিম্বী, বিপ্রতিসম ও পরিযায়ী।

স্বজ্ঞাগতভাবে এই জগৎগুলিকে জ্ঞানের অবস্থা বা ``প্রমাণসংক্রান্ত
পরিস্থিতি'' হিসেবে ভাবা যায়। $w$ থেকে $w'$ অভিগম্য হওয়ার অর্থ,
আমাদের জানা তথ্য অনুযায়ী $w'$ জ্ঞানের একটি সম্ভাব্য (ভবিষ্যৎ)
অবস্থা; অর্থাৎ $w$-তে যা জানা আছে, তার সঙ্গে $w'$ সঙ্গতিপূর্ণ।
কোনো বচন একবার জানা হলে পরে তা অজানা হয়ে যায় না: $w$-তে $!A$
জানা থাকলে এবং $Rww'$ হলে $w'$-তেও $!A$ জানা থাকে। তাই
অভিগম্যতা-সম্পর্কের সাপেক্ষে ``জ্ঞান'' একঘেয়ে।

জ্ঞানতাত্ত্বিক যুক্তিবিদ্যার মতো ``$!A$ জানা'' বলতে যদি ``সব
জ্ঞানগত বিকল্পে সত্য'' বোঝাই, তবে $w$-তে $!A \land !B$ জানা
থাকার অর্থ সব জ্ঞানগত বিকল্পে $!A$ ও $!B$ উভয়ই জানা। কিন্তু জ্ঞান
একঘেয়ে এবং $R$ প্রতিবিম্বী বলে, $w$-তে $!A \land !B$ জানা
থাকবে যদি এবং কেবল যদি $w$-তে $!A$ ও $!B$ দুটিই জানা থাকে। একই
কারণে $w$-তে $!A \lor !B$ জানা থাকবে যদি এবং কেবল যদি অন্তত
একটি জানা থাকে। ফলে $\land$ ও $\lor$-এর সত্যতার শর্ত ধ্রুপদি
যুক্তিবিদ্যার শর্তের সঙ্গে মেলে।

কিন্তু শর্তসূচকের সত্যতার শর্ত ধ্রুপদি যুক্তিবিদ্যা থেকে আলাদা।
$w$-তে $!A \lif !B$ জানা থাকবে যদি এবং কেবল যদি $Rww'$-যুক্ত
কোনো $w'$-তেই $!B$ না জেনে $!A$ জানা না থাকে। এটি $w$-তে
$!A$ অজানা অথবা $!B$ জানা থাকার শর্তের সমান নয়। কারণ $w$-তে
$!A$ ও $!B$ কোনোটিই জানা না থাকলেও এমন ভবিষ্যৎ জ্ঞানাবস্থা~$w'$
থাকতে পারে, যেখানে $Rww'$ এবং $!B$ না জেনেও $!A$ জানা যায়।

আমরা $\lnot !A$ জানি কেবল তখনই, যখন ভবিষ্যতের কোনো সম্ভাব্য
জ্ঞানাবস্থায় $!A$ জানা সম্ভব নয়। ধারণাটি হলো: $!A$ জানা সম্ভব
হলে ভবিষ্যতের কোনো সম্ভাব্য জ্ঞানাবস্থায় $!A$ জানা যাবে। যেহেতু
$\lfalse$ জানা যায় না, সেই ভবিষ্যৎ অবস্থায় $!A$ জানা থাকবে,
কিন্তু $\lfalse$ জানা থাকবে না।

এই ব্যাখ্যায় বর্জিত-মধ্যম নীতি ব্যর্থ হয়। এমন কোনো কোনো $!A$ আছে
যা আমরা এখনো জানি না, তবে ভবিষ্যতে জানতে পারি। এ রকম
!!a{formula}~$!A$-এর ক্ষেত্রে $!A$ ও $\lnot !A$ দুটিই অজানা;
সুতরাং $!A \lor \lnot !A$ জানা নেই। কিন্তু, যেমন,
$\lnot(!A \land \lnot !A)$ আমরা জানি। কারণ ভবিষ্যতের কোনো
অবস্থায় $!A$ ও $\lnot !A$ দুটিই জানা সম্ভব নয়; আর $!A$ বা
$\lnot !A$ জানা আছে কি না, তার ওপর এই জ্ঞান নির্ভর করে না।

!!^{relational model}s-ই স্বজ্ঞাবাদী যুক্তিবিদ্যার একমাত্র
অর্থবিজ্ঞান নয়। টপোলজিক্যাল অর্থবিজ্ঞান আরেকটি পদ্ধতি: এখানে
বচনগুলিকে টপোলজিক্যাল স্থানের উন্মুক্ত সেট হিসেবে এবং সংযোজকগুলিকে
সেই সেটগুলির ওপর ক্রিয়া হিসেবে ব্যাখ্যা করা হয় (যেমন,
$\land$-এর সঙ্গে ছেদের মিল আছে)।

\end{document}
